\section{Problem Statement}\label{sec:intro:problem}

Passive \gls{esm} produces a time-ordered stream of \glspl{pdw} in which pulses from several emitters are interleaved and each emitter may change mode over time (\cref{sec:intro:background}).
Operators need both a grouping by emitter and a grouping by mode from the same stream.
Classical pipelines usually deinterleave first and then apply separate mode heuristics or classifiers.
That split does not produce one representation whose geometry shows both partitions at once, and it scales poorly when waveforms are agile and pulse trains overlap.

This thesis treats the two partitions as a joint learning problem.
Training data come from a scenario simulator (\cref{sec:method:data}) in which every pulse has two labels: an emitter identity that is unique only inside its recording, and a mode drawn from a catalogue shared across recordings.
The encoder maps each window of pulses to two embedding sequences.
A clustering step on each sequence recovers the partitions, without assigning names from a fixed catalogue.
Emitter clustering is recording-local and does not assume a known number of emitters.
Mode labels are used in training to pull same-mode pulses together; at inference, mode clusters are also formed per window.
The experiments measure agreement with the simulator partitions under mixture, noise, and a fixed window length.
